
\slajdid{05-s002}
\begin{frame}[label=05-s002]{Poređenje - Komparatori}
\gusttekst
Postoje standardne kombinacione mreže binarni komparatori koji kao ulaze primaju dva binarna broja $A$ i $B$, a na izlazima daju rezultat poređenja $A<B$, $A=B$ i $A>B$. Isto kao i kod kodera prioriteta moguće je funkcije pojedinih izlaza minimizirati, međutim i komercijalno raspoložive komponente a i ovde ćemo prikazati komponentu kod koje je lako pratiti putanje signala, kao i lako nadograđivati. Primer je poređenje dva neoznačena dvobitna binarna broja
\begin{columns}[c,onlytextwidth]
\column{.77\textwidth}
\def\DEprosiren{0}\def\DEunit{.75cm}% Parametri: \DEprosiren=0/1; \DEunit određuje geometriju, font ostaje 8 pt.
\begin{tikzpicture}[x=\DEunit,y=\DEunit,circuit logic US,font=\fontsize{8}{9}\selectfont,
 g/.style={scale=.8,draw,and gate US,logic gate inputs=nn,minimum width=6mm,minimum height=4mm},
 o/.style={scale=.8,draw,or gate US,logic gate inputs=nn,minimum width=6mm,minimum height=4mm},
 n/.style={scale=.8,draw,nor gate US,logic gate inputs=nn,minimum width=6mm,minimum height=4mm},
 inv/.style={draw,not gate US,minimum width=4mm,minimum height=4mm}]
\foreach \bit/\ya/\yb/\ye in {1/4/2.7/3.35,0/1/-.3/.35}{
 \node[inv] (na\bit) at (.5,\ya) {};
 \node[inv] (aa\bit) at (1.5,\ya) {};
 \node[inv] (nb\bit) at (.5,\yb) {};
 \node[inv] (bb\bit) at (1.5,\yb) {};
 \draw[DEvod] (na\bit.input)--++(-.5,0) node[left]{A\bit};
 \draw[DEvod] (nb\bit.input)--++(-.5,0) node[left]{B\bit};
 \draw[DEvod] (na\bit.output)--(aa\bit.input);
 \draw[DEvod] (nb\bit.output)--(bb\bit.input);
 \coordinate (ba\bit) at ($(na\bit.output)!.5!(aa\bit.input)$);
 \coordinate (bbt\bit) at ($(nb\bit.output)!.5!(bb\bit.input)$);
 \node[g] (p\bit) at (3.3,\ya) {};
 \node[g] (q\bit) at (3.3,\yb) {};
 \node[n] (e\bit) at (5.3,\ye) {};
 \node[above,inner sep=1pt] at (p\bit.north) {A\bit>B\bit};
 \node[below,inner sep=1pt] at (q\bit.south) {A\bit<B\bit};
 \node[anchor=west,inner sep=0pt,text=DEplava] at (6.7,\ye+.30) {A\bit=B\bit};
 \draw[DEplava,densely dotted,line width=.35pt,-{Latex[length=1mm]}] (6.65,\ye+.30)--(e\bit.output);
 \draw[DEvod] (aa\bit.output)--++(.45,0)|-(p\bit.input 1);
 \draw[DEvod] (bb\bit.output)--++(.45,0)|-(q\bit.input 2);
 \draw[DEvod] (bbt\bit)--++(0,.25)--(p\bit.input 2);
 \draw[DEvod] (ba\bit)--++(0,-.25)--(q\bit.input 1);
 \fill (ba\bit) circle (.8pt);\fill (bbt\bit) circle (.8pt);
 \draw[DEvod] (p\bit.output)--(4.3,\ya)|-(e\bit.input 1);
 \draw[DEvod] (q\bit.output)--(4.5,\yb)|-(e\bit.input 2);
}
\node[g] (gp) at (9.0,1.75) {};
\node[g] (gq) at (9.0,.15) {};
\draw[DEvod] (e1.output)--(6.5,3.35)--(6.5,.35)|-(gq.input 1);
\draw[DEvod] (6.5,3.35)|-(gp.input 1);
\fill (6.5,3.35) circle (.8pt);
\draw[DEvod] (p0.output)--(7,1)|-(gp.input 2);
\draw[DEvod] (q0.output)--(7,-.3)|-(gq.input 2);
\node[o] (op) at (12.0,2.25) {};
\node[o] (oq) at (12.0,.65) {};
\node[g] (oe) at (12.0,-.95) {};
\draw[DEvod] (p1.output)--(4.05,4)--(10.7,4)|-(op.input 1);
\draw[DEvod] (q1.output)--(4.05,2.7)--(10.5,2.7)|-(oq.input 1);
\draw[DEvod] (gp.output)--(11.0,1.75)|-(op.input 2);
\draw[DEvod] (gq.output)--(11.0,.15)|-(oq.input 2);
\draw[DEvod] (6.5,.35)--(6.5,-.8)|-(oe.input 1);
\fill (6.5,.35) circle (.8pt);
\draw[DEvod] (e0.output)--(6.1,.35)--(6.1,-1.2)|-(oe.input 2);
\ifnum\DEprosiren=1
 \draw[DEvod] (-.5,5.9) node[left]{inA>B}--(15.8,5.9);
 \draw[DEvod] (-.5,5.5) node[left]{inA<B}--(15.5,5.5);
 \draw[DEvod] (-.5,5.1) node[left]{inA=B}--(13.1,5.1)--(13.1,-.7);
 \foreach \name/\yy in {p/2.25,q/.65,e/-.95}{
  \node[g] (c\name) at (14.3,\yy) {};
  \draw[DEvod] (o\name.output)--++(.8,0)|-(c\name.input 2);
  \draw[DEvod] (13.1,\yy+.25)|-(c\name.input 1);
  \fill (13.1,\yy+.25) circle (.8pt);
 }
 \node[o] (fp) at (16.6,2.25) {};
 \node[o] (fq) at (16.6,.65) {};
 \draw[DEvod] (15.8,5.9)|-(fp.input 1);
 \draw[DEvod] (15.5,5.5)|-(fq.input 1);
 \draw[DEvod] (cp.output)--++(.7,0)|-(fp.input 2);
 \draw[DEvod] (cq.output)--++(.7,0)|-(fq.input 2);
 \draw[DEvod] (fp.output)--++(.35,0) node[right]{A>B};
 \draw[DEvod] (fq.output)--++(.35,0) node[right]{A<B};
 \draw[DEvod] (ce.output)--(17.3,-.95) node[right]{A=B};
\else\ifnum\DEprosiren=2
 \draw[DEvod] (oe.output)--(13.1,-.95)--(13.1,2.5);
 \foreach \name/\yy in {p/2.25,q/.65,e/-.95}{
  \node[g] (c\name) at (14.3,\yy) {};
  \draw[DEvod] (13.1,\yy+.25)|-(c\name.input 1);
  \fill (13.1,\yy+.25) circle (.8pt);
 }
 \draw[DEvod] (-.5,-1.8) node[left]{inA=B}--(13.2,-1.8)|-(ce.input 2);
 \draw[DEvod] (-.5,-2.3) node[left]{inA>B}--(13.4,-2.3)|-(cp.input 2);
 \draw[DEvod] (-.5,-2.8) node[left]{inA<B}--(13.6,-2.8)|-(cq.input 2);
 \node[o] (fp) at (16.6,2.25) {};
 \node[o] (fq) at (16.6,.65) {};
 \draw[DEvod] (op.output)--(12.65,2.25)--(12.65,3.2)--(15.3,3.2)|-(fp.input 1);
 \draw[DEvod] (oq.output)--(12.8,.65)--(12.8,1.6)--(15.1,1.6)|-(fq.input 1);
 \draw[DEvod] (cp.output)--++(.7,0)|-(fp.input 2);
 \draw[DEvod] (cq.output)--++(.7,0)|-(fq.input 2);
 \draw[DEvod] (fp.output)--++(.35,0) node[right]{A>B};
 \draw[DEvod] (fq.output)--++(.35,0) node[right]{A<B};
 \draw[DEvod] (ce.output)--(17.3,-.95) node[right]{A=B};
\else
 \draw[DEvod] (op.output)--++(.35,0) node[right]{A>B};
 \draw[DEvod] (oq.output)--++(.35,0) node[right]{A<B};
 \draw[DEvod] (oe.output)--++(.35,0) node[right]{A=B};
\fi\fi
\end{tikzpicture}

\column{.21\textwidth}\centering
\begin{tikzpicture}[font=\fontsize{8}{9}\selectfont]
\draw[DEvod] (0,0) rectangle (2.6,1.7);
\foreach \label/\x in {A1/.3,A0/.8,B1/1.65,B0/2.2}{
\node[anchor=north,inner sep=1pt] at (\x,1.7) {\label};\draw[DEvod] (\x,1.7)--(\x,2.0);}
\node[anchor=west] at (0,.8) {komparator};
\foreach \label/\y in {A>B/1.15,A<B/.72,A=B/.29}{\node[anchor=east,inner sep=1pt] at (2.6,\y){\label};\draw[DEvod] (2.6,\y)--(2.8,\y);}
\end{tikzpicture}
\end{columns}
\end{frame}
